\documentclass[12pt,a4paper,oneside]{report}
\usepackage{fontspec}
\IfFontExistsTF{Times New Roman}{\setmainfont{Times New Roman}}{\setmainfont{FreeSerif}}
\usepackage[english,vietnamese]{babel}
\usepackage{geometry}
\geometry{left=3.0cm,right=2.5cm,top=2.5cm,bottom=2.5cm}
\usepackage{setspace}
\onehalfspacing
\usepackage{amsmath,amssymb}
\usepackage{booktabs,tabularx,longtable,array}
\usepackage{enumitem}
\usepackage{float}
\usepackage{hyperref}
\hypersetup{hidelinks,pdftitle={Nghiên cứu trợ lý mua sắm ứng dụng RAG và AI Agent}}
\newcommand{\srcpath}[1]{\texttt{\detokenize{#1}}}
\renewcommand{\bibname}{Tài liệu tham khảo}
\begin{document}

\begin{titlepage}
\centering
{\large  NGHIÊN CỨU KHOA HỌC\par}
\vspace{3.4cm}
{\bfseries\LARGE NGHIÊN CỨU TRỢ LÝ MUA SẮM ỨNG DỤNG RAG VÀ AI AGENT NHẰM TỰ ĐỘNG HÓA QUY TRÌNH ĐẶT HÀNG TRONG THƯƠNG MẠI ĐIỆN TỬ\par}
\vspace{2cm}
{\large Lĩnh vực: Công nghệ thông tin và thương mại điện tử\par}
\vfill
{\large Bản thảo nghiên cứu, tháng 10 năm 2026\par}
\end{titlepage}

\pagenumbering{roman}
\chapter*{Tóm tắt}
\addcontentsline{toc}{chapter}{Tóm tắt}
Báo cáo nghiên cứu cách kết hợp tạo sinh có tăng cường truy xuất (RAG) với AI Agent để hỗ trợ người mua từ mô tả nhu cầu đến tìm sản phẩm, hỏi chính sách, chuẩn bị giỏ hàng và tra cứu đơn. Giải pháp phân tuyến yêu cầu theo miền dữ liệu: mô tả sản phẩm và tài liệu chính sách được truy xuất từ chỉ mục; giá, tồn kho, giỏ hàng và đơn hàng được đọc từ dịch vụ nghiệp vụ. Các hành động thay đổi dữ liệu cần xác nhận rõ ràng, còn dữ liệu đơn hàng phải gắn với tài khoản đã xác thực. Báo cáo mô tả kiến trúc nguyên mẫu trong dự án, đề xuất bộ kịch bản và phép đo để so sánh với quy trình thủ công và chatbot chỉ trả lời thông tin. Hiện chưa có kết quả thực nghiệm trên bộ đánh giá độc lập; vì vậy các chỉ số trong báo cáo là tiêu chí và quy trình đo, không phải kết quả đạt được.

\noindent\textbf{Từ khóa:} RAG; AI Agent; trợ lý mua sắm; tìm kiếm đa phương thức; thương mại điện tử; xác nhận của người dùng.

\tableofcontents
\clearpage\pagenumbering{arabic}

\chapter{Mở đầu}
\section{Lý do chọn đề tài}
Trong cách mua sắm thông thường trên sàn thương mại điện tử, người mua tự nhập từ khóa, duyệt danh mục và sử dụng bộ lọc để thu hẹp kết quả. Sau đó, họ mở từng trang sản phẩm để đối chiếu mô tả, giá, đánh giá và tình trạng còn hàng; lựa chọn phân loại, số lượng rồi thêm vào giỏ. Trước khi đặt hàng, người mua tiếp tục kiểm tra giỏ, địa chỉ nhận hàng, ưu đãi và phương thức thanh toán. Khi phát sinh nhu cầu theo dõi hoặc đổi trả, họ phải tìm đến thông tin đơn hàng và chính sách liên quan. Các bước này có thể được thực hiện trên cùng một sàn, nhưng người mua vẫn phải tự kết nối thông tin giữa nhiều trang và thao tác. Với những yêu cầu chưa được xác định rõ, chẳng hạn tìm một sản phẩm theo mô tả hoặc hỏi về đơn hàng theo tên sản phẩm, cách tương tác này đòi hỏi thêm thời gian để tìm kiếm và đối chiếu. Thực tế đó đặt ra nhu cầu nghiên cứu một trợ lý có thể tiếp nhận yêu cầu bằng ngôn ngữ tự nhiên và hỗ trợ người mua xuyên suốt quá trình mua hàng.

Các phương thức tìm kiếm và tư vấn thông thường chủ yếu giúp người dùng tra cứu thông tin. Tuy nhiên, thông tin trong thương mại điện tử có nhiều dạng và thay đổi theo thời gian: mô tả sản phẩm và chính sách thường là văn bản, trong khi giá, tồn kho, giỏ hàng và trạng thái đơn hàng được quản lý bởi các dịch vụ nghiệp vụ. Nếu chỉ dựa vào mô hình ngôn ngữ để tạo câu trả lời, hệ thống có thể đưa ra thông tin không có căn cứ hoặc không còn phù hợp với dữ liệu hiện tại. Kỹ thuật Retrieval-Augmented Generation (RAG) cho phép truy xuất dữ liệu liên quan trước khi tạo câu trả lời, qua đó hỗ trợ việc tư vấn dựa trên thông tin sản phẩm và chính sách của hệ thống. Dù vậy, một câu trả lời đúng về sản phẩm vẫn chưa giải quyết được toàn bộ nhu cầu nếu người mua muốn thêm sản phẩm vào giỏ hoặc kiểm tra một đơn hàng cụ thể.

Việc kết hợp RAG với AI Agent tạo cơ sở để mở rộng trợ lý mua sắm từ chức năng hỏi đáp sang hỗ trợ tác vụ. Tùy theo yêu cầu, hệ thống có thể chuyển đến nhánh tìm kiếm sản phẩm, trả lời chính sách, tra cứu đơn hàng hoặc gọi công cụ để tương tác với các dịch vụ phía sau. Chẳng hạn, sau khi xác định đúng sản phẩm và phân loại, trợ lý có thể chuẩn bị thao tác thêm vào giỏ hàng; khi người dùng hỏi về tình trạng giao hàng, trợ lý cần đọc dữ liệu đơn thuộc tài khoản đã xác thực thay vì suy đoán từ lịch sử hội thoại. Những thao tác làm thay đổi dữ liệu cần có bước xác nhận rõ ràng. Đặc biệt, việc kiểm tra lần cuối, xác nhận đặt hàng và thanh toán vẫn do người mua thực hiện. Phạm vi tự động hóa của đề tài vì thế tập trung vào các bước hỗ trợ trước khi hoàn tất giao dịch.

Bên cạnh khả năng thực hiện tác vụ, bài toán nghiên cứu còn đặt ra yêu cầu về độ chính xác, an toàn và chi phí vận hành. Trợ lý phải phân biệt được câu hỏi tìm sản phẩm với câu hỏi về đơn đã mua; xử lý trường hợp một tên sản phẩm khớp nhiều đơn; kiểm tra giá và tồn kho trước khi đưa ra gợi ý; đồng thời bảo đảm thông tin đơn hàng chỉ được truy cập theo danh tính của người mua. Việc điều phối các nhánh xử lý cũng cần hợp lý để hạn chế những lần gọi mô hình hoặc truy xuất dữ liệu không cần thiết, đặc biệt khi sử dụng API có giới hạn về số lượt gọi và chi phí. Đây là các vấn đề có thể được nghiên cứu, thiết kế giải pháp và đánh giá bằng những tình huống sử dụng cụ thể.

Từ những lý do trên, đề tài “Nghiên cứu trợ lý mua sắm ứng dụng RAG và AI Agent nhằm tự động hóa quy trình đặt hàng trong thương mại điện tử” được lựa chọn nhằm nghiên cứu cách kết hợp truy xuất thông tin với năng lực điều phối và gọi công cụ trong một hệ thống thống nhất. Kết quả nghiên cứu hướng đến việc xây dựng trợ lý có khả năng trả lời dựa trên dữ liệu, hỗ trợ các thao tác mua sắm phù hợp với ngữ cảnh và giữ người mua ở vị trí quyết định đối với giao dịch cuối cùng. Đề tài đồng thời tạo điều kiện đánh giá các tiêu chí như độ chính xác của câu trả lời, khả năng nhận diện yêu cầu, mức độ hoàn thành tác vụ, thời gian phản hồi và chi phí sử dụng mô hình.

\section{Mục tiêu, câu hỏi nghiên cứu và phạm vi}
Mục tiêu là thiết kế và đánh giá một trợ lý hội thoại có thể (i) tìm và giải thích thông tin có nguồn; (ii) chọn đúng dịch vụ nghiệp vụ khi cần dữ liệu hiện hành; (iii) hỗ trợ thao tác giỏ hàng trong ranh giới xác nhận của người mua. Ba câu hỏi nghiên cứu là:
\begin{enumerate}[label=RQ\arabic*.]
\item Việc kết hợp truy xuất văn bản, hình ảnh và kiểm tra dữ liệu nghiệp vụ ảnh hưởng thế nào đến độ phù hợp của kết quả tìm sản phẩm và tính đúng đắn của câu trả lời?
\item Điều phối theo nhánh tác vụ có cải thiện tỷ lệ hoàn thành và giảm số thao tác của người mua so với chatbot RAG chỉ cung cấp thông tin không?
\item Các điểm xác nhận và kiểm tra quyền truy cập có ngăn được thao tác sai, truy cập đơn hàng trái quyền và sử dụng giá/tồn kho cũ không?
\end{enumerate}
Phạm vi đánh giá chính dừng ở tư vấn, chuẩn bị giỏ hàng, chuyển bước thanh toán và tra cứu đơn. Nguyên mẫu mã nguồn có thêm luồng tạo đơn sau một lượt xác nhận cuối tách biệt; luồng này được mô tả như khả năng triển khai nhưng không được tính là bằng chứng agent tự thanh toán. Thanh toán ví cần xác thực bổ sung từ người mua; chuyển khoản chỉ tạo đơn ở trạng thái chờ cho đến khi có xác nhận từ luồng thanh toán.

\chapter{Tổng quan và khoảng trống nghiên cứu}
\section{Khảo sát các nền tảng tại Việt Nam}
Phạm vi khảo sát. Theo báo cáo thị trường năm 2025 của Metric, bốn sàn bán lẻ trực tuyến lớn được theo dõi là Shopee, TikTok Shop, Lazada và Tiki, với tổng doanh số giao dịch khoảng 429,7 nghìn tỷ đồng. Để có năm trường hợp đối chiếu, nghiên cứu bổ sung Chợ Tốt, một nền tảng giao dịch trực tuyến có mô hình thiên về tin đăng và kết nối người mua với người bán. Chợ Tốt được xem là trường hợp so sánh về trải nghiệm tìm kiếm và giao tiếp, không được xếp ngang với bốn sàn trên về quy mô doanh số hoặc quy trình giao dịch. Sendo và Sendo Farm không được đưa vào nhóm nền tảng đang hoạt động do đã dừng vận hành.. \cite{metric2025,vneconomy2026}

Shopee. Shopee cung cấp quy trình tìm kiếm, lựa chọn sản phẩm, đặt hàng và theo dõi đơn trên ứng dụng. Đáng chú ý, người mua có thể nhập mã đơn tại mục “Trò chuyện với Shopee” để nhận thông tin trạng thái; vì vậy, không thể nhận định rằng Shopee chỉ có chức năng tìm kiếm sản phẩm hoặc không hỗ trợ tra cứu đơn qua hội thoại. Sea Limited cũng công bố các chatbot AI đang xử lý nhiều yêu cầu chăm sóc khách hàng và cho biết họ đang thử nghiệm trợ lý mua sắm AI cho người mua, với kế hoạch mở rộng trong tương lai. Công bố này chưa đủ để kết luận phạm vi triển khai cụ thể của trợ lý thử nghiệm tại Việt Nam hoặc khả năng thực hiện từng thao tác mua hàng bằng hội thoại.. \cite{shopeeorder,sea2026}

TikTok Shop. TikTok Shop kết hợp hoạt động khám phá sản phẩm với nội dung video, livestream và khu vực mua sắm. Theo quy trình dành cho người mua do nền tảng công bố, khách hàng xem sản phẩm, thêm vào giỏ, cung cấp thông tin giao hàng, chọn phương thức thanh toán và xác nhận mua hàng trên ứng dụng. Ở phía người bán tại Việt Nam, TikTok Shop đã giới thiệu Seller Assistant: công cụ có thể trả lời câu hỏi, khai thác dữ liệu vận hành và hỗ trợ một số hành động nghiệp vụ. Tuy nhiên, trợ lý dành cho người bán không thể được xem là bằng chứng rằng người mua đã có cùng một trải nghiệm trợ lý thực hiện toàn bộ chuỗi tìm kiếm, chuẩn bị giỏ hàng và tra cứu đơn bằng hội thoại.. \cite{tiktokseller,tiktokshop}

Lazada. Lazada là trường hợp cần đối chiếu kỹ nhất với đề tài. Tháng 10/2025, nền tảng công bố hệ thống gồm AI Lazzie hỗ trợ tư vấn mua sắm cùng các agent phục vụ hoàn tiền và logistics. Theo mô tả của Lazada, AI Lazzie hiểu yêu cầu hội thoại và gợi ý sản phẩm; các agent khác hỗ trợ xử lý vấn đề sau mua và thông tin vận chuyển. Như vậy, sẽ không chính xác nếu khẳng định ý tưởng kết hợp trợ lý mua sắm với nhiều agent chưa từng xuất hiện trên thị trường. Tuy nhiên, thông cáo ở cấp khu vực không cung cấp đầy đủ bằng chứng để xác định mọi chức năng đã được triển khai giống nhau tại Việt Nam, cũng như các bước nào agent có thể thực hiện thay người mua trong quy trình đặt hàng.. \cite{lazada2025}

Tiki. Tiki cung cấp kênh Tiki Chat để người mua trao đổi với nhà bán về sản phẩm. Tài liệu hướng dẫn của Tiki mô tả các chức năng như gửi hình ảnh, thiết lập tin nhắn trả lời tự động và gửi mã giảm giá. Tài liệu cũng phân biệt phạm vi tư vấn sản phẩm của nhà bán với các vấn đề về đơn hàng, đổi trả cần được chuyển đến bộ phận chăm sóc khách hàng. Đây là bằng chứng về cách tổ chức kênh hỗ trợ được công bố, nhưng không đủ để kết luận Tiki không sử dụng AI ở bất kỳ bộ phận nào khác.. \cite{tikichat}

Chợ Tốt. Chợ Tốt cho phép người mua tìm tin đăng bằng từ khóa hoặc danh mục, lọc theo khu vực và đặc điểm sản phẩm, xem chi tiết tin rồi liên hệ người bán. Quy trình này phản ánh mô hình giao dịch dựa nhiều vào việc tìm tin và trao đổi giữa hai bên. Vì cơ chế mua bán của Chợ Tốt khác với quy trình giỏ hàng và thanh toán thống nhất của các sàn bán lẻ nói trên, trường hợp này chủ yếu giúp đối chiếu mức độ hỗ trợ người dùng từ lúc nêu nhu cầu đến lúc liên hệ và ra quyết định mua.. \cite{chototsearch}

\section{Khoảng trống nghiên cứu}
Kết quả khảo sát cho thấy các nền tảng thương mại điện tử đã ứng dụng nhiều hình thức hỗ trợ người mua. Shopee cho phép tra cứu trạng thái đơn hàng qua kênh trò chuyện và đang thử nghiệm trợ lý mua sắm AI; Lazada đã công bố AI Lazzie cùng các agent phục vụ hoàn tiền và logistics. Vì vậy, sự hiện diện của chatbot hoặc AI Agent không còn là điểm mới đủ để xác lập đóng góp cho đề tài. Từ các tài liệu công khai được khảo sát, vấn đề cần nghiên cứu sâu hơn là cách kết nối thông tin được truy xuất với hành động trên hệ thống nghiệp vụ trong một luồng hội thoại, đồng thời xác định rõ mức độ tự động hóa và quyền kiểm soát của người mua. Đây là khoảng trống trong phạm vi khảo sát và đánh giá của đề tài, không phải kết luận rằng các nền tảng hiện có thiếu những chức năng đó.. \cite{shopeeorder,sea2026,lazada2025}

Khoảng trống thứ nhất liên quan đến sự liên tục của quá trình hỗ trợ mua hàng. Một yêu cầu của người mua có thể bắt đầu bằng mô tả hoặc hình ảnh sản phẩm, tiếp diễn bằng câu hỏi về đặc điểm và chính sách, sau đó chuyển thành yêu cầu chọn phân loại, thêm vào giỏ hoặc kiểm tra đơn đã mua. Mỗi bước sử dụng một loại dữ liệu khác nhau. RAG phù hợp với việc truy xuất và tổng hợp thông tin sản phẩm, chính sách; trong khi giá, tồn kho, giỏ hàng và trạng thái đơn cần được đọc hoặc cập nhật qua các dịch vụ nghiệp vụ. Do đó, cần nghiên cứu cơ chế điều phối để trợ lý nhận biết đúng ý định, sử dụng đúng nguồn dữ liệu và duy trì ngữ cảnh khi người dùng chuyển giữa các tác vụ, thay vì xử lý từng câu hỏi như một yêu cầu độc lập.

Khoảng trống thứ hai là sự cân bằng giữa tự động hóa và an toàn giao dịch. Việc hiểu câu “mua hai cái size L” chưa đủ để hệ thống thực hiện ngay thao tác thêm vào giỏ: trợ lý còn phải xác định đúng sản phẩm, đúng phân loại, kiểm tra dữ liệu hiện hành và nhận xác nhận của người mua. Với câu hỏi về đơn hàng, hệ thống phải giới hạn dữ liệu theo tài khoản đã xác thực và yêu cầu làm rõ khi có nhiều đơn cùng khớp với tên sản phẩm. Những tình huống này đặt ra bài toán thiết kế ranh giới giữa các hành động trợ lý được phép thực hiện và các quyết định phải dành cho con người. Trong phạm vi đề tài, trợ lý có thể hỗ trợ chuẩn bị giỏ hàng và chuyển đến bước thanh toán; người mua vẫn tự kiểm tra, xác nhận đặt hàng và thực hiện thanh toán.

Khoảng trống thứ ba thuộc về đánh giá hiệu quả của giải pháp. Một trợ lý tạo được câu trả lời tự nhiên chưa chắc đã tìm đúng sản phẩm, chọn đúng đơn hàng hoặc thực hiện đúng thao tác mà người mua yêu cầu. Vì vậy, đề tài cần đánh giá trên các tình huống hội thoại gồm nhiều bước, với những chỉ số cụ thể như độ chính xác khi nhận diện ý định, mức độ phù hợp của thông tin truy xuất, tỷ lệ hoàn thành tác vụ, số bước người mua phải tự thực hiện, thời gian phản hồi và lượng token sử dụng. Các trường hợp mơ hồ, dữ liệu thay đổi hoặc yêu cầu vượt quyền truy cập cũng cần được đưa vào đánh giá để kiểm tra khả năng từ chối hoặc yêu cầu làm rõ của hệ thống.

Từ những khoảng trống trên, đề tài tập trung nghiên cứu một kiến trúc trợ lý mua sắm kết hợp RAG và AI Agent, trong đó tác tử điều phối yêu cầu, truy xuất thông tin khi cần và gọi công cụ để hỗ trợ một phần quy trình mua hàng. Đóng góp dự kiến của nghiên cứu là một mô hình triển khai và bộ tiêu chí đánh giá có thể kiểm chứng, làm rõ mức độ mà hội thoại thông minh có thể giảm thao tác cho người mua mà vẫn duy trì tính chính xác, an toàn và quyền quyết định của con người.

\section{Đề xuất giải pháp}
Đề tài đề xuất xây dựng AI Agent nhằm tự động hóa các bước hỗ trợ trong quy trình mua hàng trên sàn thương mại điện tử. Từ yêu cầu bằng ngôn ngữ tự nhiên hoặc hình ảnh, trợ lý xác định nhu cầu của người mua, sử dụng RAG để truy xuất thông tin sản phẩm và chính sách, sau đó gọi các công cụ kết nối với hệ thống nghiệp vụ để kiểm tra giá, tồn kho, phân loại và trạng thái giỏ hàng. Khi người mua nêu rõ sản phẩm, số lượng và phân loại mong muốn, agent chuẩn bị thao tác thêm vào giỏ và chỉ thực hiện sau khi nhận được xác nhận. Agent cũng có thể tra cứu các đơn hàng thuộc tài khoản đã xác thực để hỗ trợ người mua sau khi đặt hàng. Quy trình tự động hóa kết thúc ở bước chuẩn bị và chuyển sang trang thanh toán; người mua tự kiểm tra, xác nhận đặt hàng và thanh toán. Cách tiếp cận này cho phép đánh giá AI Agent theo khả năng hoàn thành tác vụ mua sắm, thay vì chỉ đánh giá chất lượng câu trả lời của chatbot.

\chapter{Cơ sở khoa học}
\section{Mô hình ngôn ngữ lớn}
Mô hình ngôn ngữ lớn (Large Language Model – LLM) là mô hình trí tuệ nhân tạo có khả năng xử lý và tạo văn bản dựa trên ngữ cảnh đầu vào. Trong trợ lý mua sắm, LLM giúp diễn giải những yêu cầu được diễn đạt bằng ngôn ngữ tự nhiên, kể cả khi người dùng không nêu đầy đủ thông tin. Ví dụ, sau khi trao đổi về một mẫu áo, yêu cầu “thêm hai cái size L” cần được hiểu là thao tác với mẫu áo vừa được nhắc đến.

Tuy nhiên, LLM không tự biết chính xác giá bán, số lượng tồn kho, nội dung giỏ hàng hay trạng thái đơn hàng tại thời điểm người dùng đặt câu hỏi. Các thông tin này phải được lấy từ nguồn dữ liệu của hệ thống. Đây là cơ sở để đề tài kết hợp LLM với hai thành phần bổ trợ: RAG để tìm kiếm tri thức phù hợp và AI Agent để lựa chọn, thực hiện các thao tác thông qua công cụ.

\section{Biểu diễn vector và tìm kiếm ngữ nghĩa}
Biểu diễn vector, hay embedding, là quá trình chuyển văn bản hoặc hình ảnh thành dãy số thể hiện một phần nội dung và đặc trưng của chúng. Những đối tượng có nội dung gần nhau có thể được tìm thấy bằng cách so sánh độ tương đồng giữa các vector. Nhờ đó, hệ thống có thể tìm sản phẩm phù hợp với câu mô tả “áo thun nam màu đen, thoáng mát” dù tên sản phẩm không chứa đúng toàn bộ cụm từ mà người dùng nhập. Nghiên cứu về truy xuất đoạn văn bản bằng biểu diễn vector cho thấy đây là một hướng tiếp cận hữu ích đối với các bài toán tìm kiếm theo nội dung ngữ nghĩa (Karpukhin và cộng sự, 2020). \cite{karpukhin2020}

Trong đề tài, tìm kiếm ngữ nghĩa đóng vai trò xác định những sản phẩm hoặc tài liệu chính sách có khả năng liên quan đến yêu cầu. Kết quả tìm kiếm vẫn cần được kiểm tra trước khi sử dụng: một sản phẩm tương đồng về mô tả chưa chắc có đúng màu, kích cỡ hoặc còn hàng. Vì vậy, dữ liệu truy xuất được dùng để chọn ứng viên và tạo câu trả lời, còn những thuộc tính có thể thay đổi như giá và tồn kho phải được đối chiếu với dữ liệu hiện hành từ hệ thống.

\section{Tạo sinh có tăng cường truy xuất}
RAG (Retrieval-Augmented Generation) là phương pháp kết hợp mô hình tạo sinh với cơ chế truy xuất thông tin từ nguồn dữ liệu bên ngoài. Thay vì chỉ dựa vào thông tin đã học trong quá trình huấn luyện, hệ thống tìm các nội dung liên quan đến câu hỏi, đưa chúng vào ngữ cảnh đầu vào, rồi mới tạo câu trả lời. Công trình của Lewis và cộng sự đặt nền tảng cho cách kết hợp tri thức lưu trong mô hình với tri thức được truy xuất từ một kho dữ liệu riêng (Lewis và cộng sự, 2020). \cite{lewis2020}

Đối với trợ lý mua sắm, RAG có thể hỗ trợ tìm sản phẩm theo nhu cầu, giải thích đặc điểm sản phẩm và trả lời các câu hỏi về chính sách đổi trả hoặc giao hàng. Quy trình cơ bản gồm chuẩn bị dữ liệu, tạo biểu diễn để tìm kiếm, truy xuất nội dung phù hợp với câu hỏi và tạo câu trả lời dựa trên nội dung đã tìm được. Chất lượng câu trả lời phụ thuộc không chỉ vào LLM mà còn vào độ đầy đủ, tính cập nhật và mức độ phù hợp của dữ liệu được truy xuất.

RAG chủ yếu giải quyết bài toán tìm và sử dụng thông tin. Bản thân việc truy xuất một mô tả sản phẩm không làm thay đổi giỏ hàng, không kiểm tra được trạng thái đơn hàng cá nhân và không thay thế cho thao tác nghiệp vụ trên hệ thống. Sự phân biệt này là cơ sở để đề tài sử dụng thêm AI Agent.

\section{Tìm kiếm sản phẩm bằng hình ảnh}
Ngoài văn bản, người mua có thể gửi ảnh của món hàng muốn tìm. Để hỗ trợ trường hợp này, hệ thống cần biểu diễn hình ảnh theo cách cho phép so sánh với hình ảnh sản phẩm trong danh mục. CLIP là một mô hình tiêu biểu học mối liên hệ giữa hình ảnh và văn bản, tạo cơ sở cho các tác vụ tìm kiếm và đối chiếu nội dung đa phương thức (Radford và cộng sự, 2021). \cite{radford2021}

Trong phạm vi đề tài, tìm kiếm bằng ảnh được xem là một phương thức xác định sản phẩm ứng viên. Kết quả tương đồng về hình ảnh cần tiếp tục được đối chiếu với tên sản phẩm, phân loại, giá và các thuộc tính khác trước khi tư vấn hoặc thực hiện thao tác mua sắm. Một hình ảnh giống mẫu mã không đủ để khẳng định hai sản phẩm hoàn toàn tương đương.

\section{AI Agent và cơ chế sử dụng công cụ}
AI Agent trong đề tài được hiểu là thành phần tiếp nhận mục tiêu của người dùng, xác định thông tin hoặc thao tác cần thiết, lựa chọn công cụ phù hợp và sử dụng kết quả trả về để tiếp tục xử lý yêu cầu. Công cụ có thể là chức năng tìm kiếm sản phẩm, đọc giỏ hàng, kiểm tra phân loại, cập nhật giỏ hàng hoặc tra cứu đơn hàng. Các nghiên cứu về kết hợp suy luận và hành động, chẳng hạn ReAct, cho thấy mô hình ngôn ngữ có thể tương tác với nguồn thông tin và môi trường bên ngoài thông qua các hành động cụ thể (Yao và cộng sự, 2023). \cite{yao2023}

Điểm cần phân biệt là LLM đề xuất hoặc lựa chọn lời gọi công cụ, còn ứng dụng và backend mới thực thi thao tác. Backend phải kiểm tra danh tính người dùng, quyền truy cập, dữ liệu đầu vào và điều kiện nghiệp vụ trước khi trả kết quả. Chẳng hạn, câu hỏi “đơn hàng của tôi đang ở đâu” phải được giải quyết bằng dữ liệu đơn hàng thuộc tài khoản đã xác thực, thay vì để mô hình suy đoán từ lịch sử hội thoại.

Một AI Agent không nhất thiết phải liên tục lập kế hoạch lại sau mỗi bước. Trong đề tài, yêu cầu được phân loại và chuyển đến nhánh xử lý phù hợp; mỗi nhánh chỉ gọi những công cụ cần thiết cho nhiệm vụ của mình. Cách tổ chức này giúp giới hạn số lần gọi mô hình và công cụ, đồng thời tạo điều kiện kiểm soát thời gian phản hồi và chi phí xử lý. Tài liệu LangGraph cũng phân biệt quy trình có tuyến xử lý được xác định trước với agent có khả năng tự quyết định việc dùng công cụ theo tình huống (LangChain, không ghi năm). \cite{langgraphworkflow}

\section{Điều phối nhánh xử lý và quản lý ngữ cảnh}
Một yêu cầu mua sắm có thể liên quan đến sản phẩm, chính sách hoặc đơn hàng. Vì nguồn dữ liệu và quyền truy cập của các nhóm yêu cầu này khác nhau, hệ thống cần cơ chế điều phối để chuyển câu hỏi đến đúng nhánh xử lý. Nhánh sản phẩm sử dụng thông tin danh mục và tìm kiếm ngữ nghĩa; nhánh chính sách sử dụng tài liệu quy định; nhánh đơn hàng truy vấn dữ liệu nghiệp vụ của người dùng đã xác thực. Việc tách nhánh giúp hạn chế tình trạng lấy kết quả tìm kiếm sản phẩm để trả lời một câu hỏi về đơn hàng, hoặc dùng thông tin chung để suy đoán dữ liệu cá nhân.

Bên cạnh việc chọn nhánh, hệ thống cần lưu giữ ngữ cảnh hội thoại ở mức cần thiết. Ngữ cảnh cho phép hiểu các câu tiếp nối như “mua hai cái size L”, “đơn đó đến đâu rồi” hoặc “có đổi được không”. Tuy nhiên, ngữ cảnh hội thoại chỉ giúp xác định đối tượng mà người dùng đang nhắc đến; nó không thay thế dữ liệu hiện hành trong cơ sở dữ liệu. Khi người dùng hỏi về đơn hàng, hệ thống vẫn phải truy vấn và kiểm tra đúng đơn hàng thuộc tài khoản của họ.

\section{Kiểm soát của con người}
Tự động hóa quy trình mua hàng trong đề tài được hiểu là giảm các thao tác thủ công ở giai đoạn tìm kiếm, tư vấn, chọn phân loại, chuẩn bị giỏ hàng và tra cứu thông tin sau mua. Những hành động làm thay đổi dữ liệu, như thêm sản phẩm vào giỏ, cần được thực hiện theo ý định rõ ràng của người dùng và có bước xác nhận phù hợp. Đối với bước đặt hàng và thanh toán, người dùng trực tiếp kiểm tra thông tin, xác nhận đơn và hoàn tất thanh toán. Cơ chế để con người xem xét và quyết định tại các điểm quan trọng cũng phù hợp với nguyên tắc human-in-the-loop trong thiết kế hệ thống có sử dụng agent (LangChain, không ghi năm). \cite{langgraphinterrupt}

Từ các khái niệm trên, đề tài xác lập vai trò bổ trợ của hai thành phần chính: RAG cung cấp thông tin liên quan để tư vấn có căn cứ; AI Agent kết nối yêu cầu bằng ngôn ngữ tự nhiên với các chức năng nghiệp vụ để hỗ trợ thực hiện quy trình mua hàng. Hiệu quả của mô hình kết hợp cần được đánh giá đồng thời theo độ chính xác của thông tin, khả năng hoàn thành tác vụ, số thao tác người dùng phải thực hiện, thời gian phản hồi, chi phí xử lý và mức độ kiểm soát đối với các hành động ảnh hưởng đến tài khoản của người mua.

\chapter{Thiết kế và hiện thực nguyên mẫu}
\section{Kiến trúc tổng thể}
Hệ thống trong kho mã nguồn gồm giao diện người mua, API Gateway, chatbot-service và các dịch vụ sản phẩm, tồn kho, giỏ hàng, đơn hàng. 
\begin{figure}[H]\centering
\fbox{\begin{minipage}{0.91\textwidth}\centering
Người mua $\longrightarrow$ API Gateway $\longrightarrow$ Router Agent\\[4pt]
$\swarrow$ Product Agent\quad $\downarrow$ Policy Agent\quad $\searrow$ Order Agent\\[4pt]
Qdrant văn bản/ảnh $+$ Product/Inventory API\quad Qdrant chính sách\quad Order API\\[4pt]
$\longrightarrow$ câu trả lời, nguồn, sản phẩm hoặc hành động chờ xác nhận
\end{minipage}}
\caption{Luồng xử lý khái quát của nguyên mẫu.}\label{fig:architecture}
\end{figure}

Router và Product Agent được tổ chức bằng LangGraph. Trạng thái hội thoại sử dụng \texttt{sessionId} làm \texttt{thread\_id} để lưu checkpoint; thông tin ngắn cần cho câu hỏi tiếp nối được giữ lại, còn ảnh gốc và dữ liệu trung gian lớn không được lưu trong checkpoint. Việc phân tuyến và lập kế hoạch tác vụ dùng quy tắc; nhánh tìm kiếm sản phẩm có một lượt truy xuất và tối đa một lần gọi mô hình để chọn, diễn đạt kết quả. Đây là đặc điểm thiết kế của nguyên mẫu, cần được đo thực tế về độ trễ và chi phí.

\section{Truy xuất và đối chiếu dữ liệu}
Văn bản sản phẩm và mô tả ảnh được nhúng vào collection văn bản của Qdrant; ảnh sản phẩm dùng CLIP trong collection ảnh riêng. Hai không gian vector không thể so sánh trực tiếp bằng điểm cosine, vì vậy kết quả được hợp nhất theo thứ hạng. Sau truy xuất, các ứng viên phải qua bộ lọc nghiệp vụ về trạng thái duyệt, tồn kho, giá và thuộc tính biến thể. Mô hình chỉ được chọn mã sản phẩm thuộc danh sách ứng viên đã lọc. Với câu hỏi chính sách, câu trả lời phải dựa trên đoạn tài liệu cùng địa chỉ nguồn; với câu hỏi đơn hàng, hệ thống dùng Order API và không suy đoán từ vector hay lịch sử chat.

\section{Thao tác giỏ hàng và điểm xác nhận}
Product Agent nhận diện yêu cầu thêm, sửa, xóa mục giỏ hoặc chuyển bước thanh toán. Trước thao tác ghi, hệ thống xác định sản phẩm, biến thể, số lượng và tài khoản; yêu cầu mơ hồ được chuyển thành câu hỏi làm rõ. Backend kiểm tra lại người mua, trạng thái catalog, tồn kho và giá. Yêu cầu ghi được lưu dưới dạng hành động chờ để người mua xác nhận ở lượt riêng. Mã yêu cầu có thể dùng để chống lặp khi thử lại.

Nguyên mẫu còn hỗ trợ báo giá thanh toán gắn với người mua và phiên trò chuyện, có thời hạn và được kiểm tra lại trước khi tạo đơn sau xác nhận cuối. Việc tạo đơn và hoàn tất thanh toán là hai trạng thái khác nhau: chuyển khoản có thể tạo đơn chờ thanh toán; thanh toán ví cần bước xác thực bổ sung trên giao diện. Báo cáo không quy kết khả năng này cho mọi nền tảng khảo sát và không xem việc gợi ý của mô hình là quyền tự thực thi giao dịch.

\section{Giả thuyết và rủi ro kỹ thuật}
Các giả thuyết có thể kiểm định gồm: kết hợp truy xuất với API hiện hành giảm câu trả lời dùng giá/tồn kho cũ; phân tuyến theo tác vụ giảm gọi mô hình không cần thiết; yêu cầu xác nhận riêng làm giảm hành động ngoài ý muốn. Các rủi ro còn lại gồm dữ liệu catalog thiếu hoặc sai, ảnh người mua khác ảnh catalog, phân tuyến nhầm giữa sản phẩm và đơn đã mua, giá thay đổi trong thời gian chờ, và độ trễ phụ thuộc API ngoài. Mỗi rủi ro cần có ca kiểm thử và chỉ số riêng, thay vì được suy ra từ kết quả tìm kiếm tổng quát.

\chapter{Phương pháp thực nghiệm và đánh giá}
\section{Phương pháp nghiên cứu}
Đề tài sử dụng phương pháp nghiên cứu thiết kế kết hợp thực nghiệm. Quy trình nghiên cứu gồm bốn giai đoạn: khảo sát bài toán và các giải pháp hiện có; thiết kế, xây dựng trợ lý mua sắm; thiết lập các tình huống kiểm thử; và đánh giá kết quả theo những tiêu chí định lượng, định tính. Đối tượng đánh giá là khả năng hỗ trợ người mua hoàn thành các bước trong quy trình mua hàng, từ tìm kiếm sản phẩm đến chuẩn bị giỏ hàng và tra cứu đơn sau mua. Phạm vi tự động hóa không bao gồm việc agent tự xác nhận đặt hàng hoặc thực hiện thanh toán.

Khảo sát và xác định yêu cầu. Nghiên cứu thu thập tài liệu công khai về quy trình mua hàng và tính năng hỗ trợ người mua của Shopee, TikTok Shop, Lazada, Tiki và Chợ Tốt. Việc khảo sát nhằm xác định những tác vụ phổ biến, cách các nền tảng tổ chức chức năng tìm kiếm, tư vấn và hỗ trợ đơn hàng, từ đó xây dựng yêu cầu cho hệ thống thực nghiệm. Song song với khảo sát thị trường, đề tài phân tích các nhóm yêu cầu có thể xuất hiện trong hội thoại: tìm sản phẩm bằng mô tả hoặc hình ảnh, hỏi về đặc điểm và chính sách, chọn phân loại và số lượng, thay đổi giỏ hàng, kiểm tra đơn hàng. Kết quả khảo sát được sử dụng để xác định phạm vi chức năng và thiết kế bộ kịch bản đánh giá; đề tài không suy luận rằng một nền tảng thiếu tính năng chỉ vì tính năng đó không được mô tả trong tài liệu công khai.

Thiết kế và xây dựng hệ thống. Trợ lý được tổ chức thành các nhánh xử lý theo nhiệm vụ. Nhánh sản phẩm sử dụng RAG để truy xuất những sản phẩm phù hợp với yêu cầu bằng văn bản hoặc hình ảnh; nhánh chính sách truy xuất tài liệu liên quan để tạo câu trả lời có căn cứ; nhánh đơn hàng đọc dữ liệu thuộc tài khoản đã xác thực để trả lời các câu hỏi sau mua. AI Agent điều phối yêu cầu đến nhánh phù hợp và gọi công cụ kết nối với hệ thống nghiệp vụ khi cần kiểm tra giá, tồn kho hoặc thao tác với giỏ hàng. Đối với hành động làm thay đổi giỏ hàng, hệ thống xác định sản phẩm, phân loại, số lượng và yêu cầu người mua xác nhận trước khi thực hiện. Sau bước chuẩn bị, trợ lý có thể chuyển người mua đến trang thanh toán, nơi họ tự kiểm tra và hoàn tất giao dịch. Luồng xử lý được thiết kế theo kế hoạch một lượt cho mỗi yêu cầu nhằm hạn chế các lần gọi mô hình không cần thiết.

Xây dựng bộ dữ liệu và kịch bản kiểm thử. Đề tài dự kiến xây dựng tối thiểu 100 kịch bản dựa trên danh mục sản phẩm, tài liệu chính sách và dữ liệu đơn hàng thử nghiệm của hệ thống. Các kịch bản được phân thành bốn nhóm: tìm kiếm và tư vấn sản phẩm; hội thoại nhiều lượt dẫn đến thao tác giỏ hàng; tra cứu đơn hàng; và hỏi đáp chính sách. Mỗi kịch bản ghi rõ yêu cầu đầu vào, dữ liệu được phép sử dụng, kết quả mong đợi và điều kiện được xem là hoàn thành. Bộ kiểm thử cần bao gồm cả trường hợp thông thường lẫn trường hợp khó, chẳng hạn yêu cầu thiếu phân loại, nhiều sản phẩm hoặc nhiều đơn hàng cùng khớp, giá và tồn kho thay đổi, người dùng chưa đăng nhập, hoặc yêu cầu truy cập đơn của tài khoản khác.

Thực nghiệm và đánh giá. Hệ thống đề xuất được đánh giá trên cùng một bộ kịch bản với hai cách tiếp cận đối chiếu: quy trình người mua tự tìm kiếm và thao tác qua giao diện, và chatbot RAG chỉ cung cấp thông tin mà không gọi công cụ thực hiện tác vụ. Các tiêu chí đánh giá gồm độ chính xác khi nhận diện ý định, mức độ phù hợp của kết quả truy xuất, tính đúng đắn của câu trả lời so với dữ liệu nguồn, tỷ lệ hoàn thành tác vụ và số bước người mua phải tự thực hiện. Về hiệu năng, nghiên cứu ghi nhận thời gian phản hồi, số lần gọi API mô hình và lượng token sử dụng cho mỗi yêu cầu. Đối với các thao tác nghiệp vụ, việc đánh giá còn kiểm tra tỷ lệ chọn đúng sản phẩm, phân loại, số lượng, cũng như việc hệ thống có yêu cầu xác nhận trước khi thay đổi giỏ hàng hay không. Kết quả được tổng hợp theo từng nhóm kịch bản để nhận diện tác vụ mà giải pháp hỗ trợ tốt và các tình huống vẫn cần người dùng làm rõ.

Kiểm soát tính an toàn và giới hạn nghiên cứu. Dữ liệu đơn hàng trong thực nghiệm chỉ được truy xuất theo danh tính của tài khoản đã xác thực. Những yêu cầu mơ hồ phải dẫn đến câu hỏi làm rõ thay vì tự chọn một đơn hoặc sản phẩm khi chưa đủ căn cứ. Các trường hợp truy cập trái quyền, thông tin sản phẩm không còn hợp lệ và yêu cầu thực hiện hành động khi chưa được xác nhận được kiểm tra riêng. Thực nghiệm được tiến hành trên hệ thống của đề tài; các nền tảng khảo sát chỉ cung cấp bối cảnh và cơ sở đối chiếu chức năng, không được xem là các hệ thống tham gia phép đo hiệu năng trực tiếp.

\section{Bộ dữ liệu, đường cơ sở và quy trình}
Bộ đánh giá dự kiến có ít nhất 100 kịch bản độc lập theo bốn nhóm trong bản thảo: tìm/tư vấn sản phẩm, hội thoại nhiều lượt với giỏ hàng, tra cứu đơn, và hỏi đáp chính sách. Mỗi kịch bản lưu ảnh hoặc lời hỏi đầu vào, ảnh chụp dữ liệu nguồn tại thời điểm chạy, danh sách mã sản phẩm/đơn hợp lệ, điều kiện làm rõ, hành động được phép và kết quả mong đợi. Các ca lỗi phải bao gồm truy cập trái quyền, hết hàng, biến thể thiếu, giá đổi và yêu cầu lặp. Cần tách kịch bản xây dựng prompt khỏi kịch bản đánh giá để tránh đo lại dữ liệu đã dùng khi phát triển.

Ba điều kiện đối chiếu được chạy trên cùng dữ liệu: (A) người dùng tự thao tác qua giao diện; (B) chatbot RAG chỉ trả lời thông tin; (C) trợ lý RAG kết hợp Agent và công cụ. Với điều kiện A, ghi số thao tác và thời gian đến mục tiêu; với B và C, ghi cùng mục tiêu cuối và số lượt cần người dùng can thiệp. Thứ tự kịch bản nên đảo giữa người tham gia để giảm hiệu ứng học. Các API và phiên bản chỉ mục phải cố định trong một đợt so sánh.

\section{Chỉ số đánh giá}
Gọi $N$ là số kịch bản, $y_i=1$ khi mục tiêu và mọi điều kiện an toàn của kịch bản $i$ đều đạt, và $y_i=0$ trong các trường hợp còn lại. Tỷ lệ hoàn thành tác vụ được tính bằng
\begin{equation}\mathrm{TCR}=\frac{1}{N}\sum_{i=1}^{N}y_i.\end{equation}
Với tìm kiếm sản phẩm, đo Recall@$k$ trên tập sản phẩm hợp lệ và tỷ lệ chọn đúng biến thể. Với hỏi đáp, người chấm đối chiếu từng mệnh đề thực chứng với đoạn nguồn hoặc API hiện hành, ghi tỷ lệ mệnh đề được hỗ trợ. Với phân tuyến, lập ma trận nhầm lẫn giữa nhãn dự đoán và nhãn chuẩn để tính precision, recall và macro-F1. Không dùng đánh giá tự động bằng cùng một mô hình làm bằng chứng duy nhất cho độ đúng.

Hiệu năng được ghi theo trung vị và phân vị 95 của thời gian phản hồi; chi phí gồm số lời gọi mô hình, số token vào/ra và ước tính chi phí theo giá tại thời điểm đo. Hiệu quả thao tác là số lượt nhấp hoặc nhập tay từ lúc nêu yêu cầu đến khi đạt mục tiêu. An toàn được báo cáo thành số lỗi trên từng loại kiểm thử: lộ đơn khác chủ, ghi dữ liệu khi chưa xác nhận, chọn nhầm sản phẩm/biến thể và dùng giá hoặc tồn kho cũ. Cần báo cáo mẫu số của từng nhóm và khoảng tin cậy khi so sánh tỷ lệ.

\section{Tình trạng bằng chứng và giới hạn}
Kho mã nguồn có script \srcpath{chatbot-service/src/scripts/evaluate_product.py} để kiểm tra tự truy xuất ảnh catalog và tra tên sản phẩm. Đây là phép kiểm tra đường ống kỹ thuật trên chính dữ liệu catalog, không thay thế phép đánh giá ảnh người mua hoặc bộ kịch bản độc lập. Báo cáo này chưa có bảng kết quả đo trên 100 kịch bản và chưa thể kết luận giải pháp cải thiện bao nhiêu so với hai đường cơ sở. Những phần mô tả chức năng dựa trên mã nguồn và tài liệu dự án; các giả thuyết hiệu quả chỉ được xác nhận sau thực nghiệm.

\chapter{Kết luận và hướng nghiên cứu tiếp theo}
Nghiên cứu xác lập mô hình trợ lý mua sắm dùng RAG cho thông tin sản phẩm/chính sách và AI Agent để điều phối công cụ nghiệp vụ, với kiểm tra dữ liệu hiện hành và xác nhận trước hành động nhạy cảm. Đóng góp hiện tại là kiến trúc và quy trình đánh giá có thể kiểm chứng. Bước tiếp theo là xây dựng bộ kịch bản độc lập, chạy ba điều kiện đối chiếu, công bố kết quả kèm dữ liệu và phân tích lỗi. Các hướng mở rộng gồm đánh giá ảnh chụp ngoài catalog, kiểm tra tính bền vững khi catalog thay đổi, và nghiên cứu xác nhận giao dịch thuận tiện nhưng vẫn rõ ràng cho người mua.

\begin{thebibliography}{99}
\bibitem{metric2025} Metric. (2026). \emph{Thị trường TMĐT Việt Nam năm 2025: những chuyển dịch đáng chú ý}. \url{https://metric.vn/insights/thi-truong-tmdt-viet-nam-nam-2025-nhung-chuyen-dich-dang-chu-y/}.
\bibitem{vneconomy2026} VnEconomy. (2026). \emph{Sendo Farm chính thức dừng hoạt động, khách hàng được hoàn tiền}. \url{https://vneconomy.vn/lan-thu-hai-sen-do-rut-chan-khoi-thuong-mai-dien-tu.htm}.
\bibitem{shopeeorder} Shopee Trung tâm trợ giúp. \emph{Theo dõi đơn hàng tại mục Trò chuyện với Shopee}. \url{https://help.shopee.vn/portal/4/article/79600}.
\bibitem{sea2026} Sea Limited. (2026). \emph{First Quarter 2026 Results Deck}. \url{https://cdn.sea.com/investor/1Q2026/HEEdeRDe0NyiZ1PrWIvH/2026.05.12%20Sea%20First%20Quarter%202026%20Results%20Deck.pdf}.
\bibitem{tiktokshop} TikTok Newsroom. \emph{TikTok Shop Tab: All you need to know}. \url{https://newsroom.tiktok.com/all-you-need-to-know-about-the-tiktok-shop-tab?lang=en-SG}.
\bibitem{tiktokseller} TikTok Shop Seller University. (2026). \emph{How Seller Assistant Works}. \url{https://seller-vn.tiktok.com/university/essay?knowledge_id=8109059019573009&lang=en}.
\bibitem{lazada2025} Lazada Group. (2025). \emph{Lazada's AI Agents Power a Smarter Shopping and Selling Experience Ahead of 11.11}. \url{https://lazada-com.oss-ap-southeast-1.aliyuncs.com/225-press-release.pdf}.
\bibitem{tikichat} Học viện Tiki. \emph{Hướng dẫn sử dụng tính năng Tiki Chat}. \url{https://hocvien.tiki.vn/faq/tinh-nang-tiki-chat/}.
\bibitem{chototsearch} Chợ Tốt Trợ giúp. \emph{Các bước tìm kiếm một sản phẩm}. \url{https://trogiup.chotot.com/nguoi-mua/cac-buoc-tim-kiem-mot-san-pham/}.
\bibitem{karpukhin2020} Karpukhin, V. và cộng sự. (2020). \emph{Dense Passage Retrieval for Open-Domain Question Answering}. \url{https://arxiv.org/abs/2004.04906}.
\bibitem{lewis2020} Lewis, P. và cộng sự. (2020). \emph{Retrieval-Augmented Generation for Knowledge-Intensive NLP Tasks}. \url{https://arxiv.org/abs/2005.11401}.
\bibitem{radford2021} Radford, A. và cộng sự. (2021). \emph{Learning Transferable Visual Models From Natural Language Supervision}. \url{https://arxiv.org/abs/2103.00020}.
\bibitem{yao2023} Yao, S. và cộng sự. (2023). \emph{ReAct: Synergizing Reasoning and Acting in Language Models}. \url{https://arxiv.org/abs/2210.03629}.
\bibitem{langgraphworkflow} LangChain. \emph{Workflows and agents}. \url{https://docs.langchain.com/oss/python/langgraph/workflows-agents}.
\bibitem{langgraphinterrupt} LangChain. \emph{Interrupts}. \url{https://docs.langchain.com/oss/python/langgraph/interrupts}.
\end{thebibliography}
\end{document}
